\documentclass[11pt]{article}
\usepackage[a4paper,margin=25mm]{geometry}
\usepackage{amsmath,amssymb,amsthm,mathtools,booktabs,array,longtable,xcolor}
\usepackage[T1]{fontenc}
\usepackage{lmodern}
\usepackage[colorlinks=true,linkcolor=teal,citecolor=teal,urlcolor=teal]{hyperref}
\hypersetup{pdftitle={Bicomplex Signal Geometry: Completed Source Charts, Monodromy and Spectral Stability},
 pdfauthor={Akari Hayami (Jian-Yu Huang)},pdfsubject={v13 working research revision}}
\newtheoremstyle{readable}{6pt}{6pt}{\itshape}{}{\bfseries}{.}{\newline}{}
\theoremstyle{readable}
\newtheorem{theorem}{Theorem}[section]
\newtheorem{proposition}[theorem]{Proposition}
\newtheorem{lemma}[theorem]{Lemma}
\newtheorem{definition}[theorem]{Definition}
\newtheorem{corollary}[theorem]{Corollary}
\theoremstyle{remark}\newtheorem{remark}[theorem]{Remark}
\newcommand{\R}{\mathbb R}\newcommand{\C}{\mathbb C}\newcommand{\Z}{\mathbb Z}
\newcommand{\Q}{\mathbb Q}\newcommand{\T}{\mathbb T}
\newcommand{\dd}{\,d}\newcommand{\norm}[1]{\left\lVert#1\right\rVert}
\newcommand{\Dom}{\operatorname{Dom}}\newcommand{\spec}{\operatorname{spec}}
\newcommand{\capg}{\operatorname{Cap}_{g}}
\DeclareMathOperator{\sgn}{sgn}\DeclareMathOperator{\rank}{rank}
\title{Bicomplex Signal Geometry:\\Completed Source Charts, Monodromy and Spectral Stability}
\author{Akari Hayami (Jian-Yu Huang)}
\date{1 October 2026\\\small v13 working revision; local research draft}
\begin{document}
\maketitle
\begin{abstract}
We study an explicit ruled map associated with a two-component signal model.
A completed source atlas includes its two previously omitted endpoint
singularities. The planar observation field has two interior zeros of opposite
index; continuous square-root lifting gives a sign change on one traversal and
restoration on two. The standard fold and cross-cap residuals have sublevel
exponent $3/4$. Their real links, an explicitly chosen complex $A_3$ completion,
and the actual principal-curvature cusp are treated as distinct constructions.
On a fixed Lebesgue Hilbert space we prove the Mellin integration-by-parts
identity, identify the half-line domain obstruction, and establish finite-time
Gram estimates. Distinct finite frequencies are independent for every positive
observation time. We also prove local zero capacity and explain why it is a
statement about closed forms rather than minimal operator domains.
Weighted radial estimates give compactness and a positive fixed-Dirichlet
gap on the completed source, including its endpoint corners.
An audit of the later operator claims exposes a concrete domain error:
closing the interior compact core leaves infinitely many outer-boundary
deficiency modes. We give a corrected conditional construction from a
Dirichlet realization and bounded point traces. The actual Green and twisted
Dirac parametrices require further estimates, which are stated explicitly.
The normalization defect, the phase sign line, and a chosen full-spin
transmission map retain different data and roles. No communication capacity,
particle assignment, or arithmetic weight follows from a numerical analogy.
\end{abstract}
\tableofcontents

\section{The model and its mathematical inputs}\label{sec:model}
The geometric object in this paper is the map
\begin{equation}\label{eq:S}
 S(t,u)=\bigl(c+2u(1-c),(1-u)s,uQ\bigr),\qquad
 c=\cos t,\quad s=\sin t,\quad Q=\sqrt{c(2-c)},
\end{equation}
with physical labels $D=[-\pi/2,\pi/2]\times[0,1]$.
We keep source labels distinct even when their images coincide. In particular,
$S(t,1)=S(-t,1)$ is an image relation, not a quotient defining the source.
The term bicomplex can be made precise by taking commuting units $i,j$ with
$i^2=j^2=-1$ and idempotents $e_\pm=(1\pm ij)/2$. The algebra then identifies
with $\C e_+\oplus\C e_-$. Choosing the coefficient norm
$|z_+|^2+|z_-|^2$ is additional Hilbert data; it is not a multiplicative
bicomplex modulus. None of the proofs below depends on bicomplex multiplication.

The momentum map $(z_+,z_-)\mapsto(|z_+|^2,|z_-|^2)$ and the unit-power
constraint define a sphere and its moment interval. They do not, by themselves,
define a two-dimensional source, a lift of \eqref{eq:S}, or a projection taking
that lift to $S$. We therefore take \eqref{eq:S} as an explicit model input.
A signal realization through an ambient projection would have to supply those
maps separately. This resolves a missing definition in historical v12,
rather than declaring that an unspecified map exists.

We use conjugate-linear first inner products throughout:
$\langle f,g\rangle=\int\overline f g$. A fold in this paper means a smooth
plane-to-plane germ equivalent to $(x,y)\mapsto(x,y^2)$. A cross-cap means a
smooth plane-to-three-space germ equivalent to $(x,y)\mapsto(x,xy,y^2)$.
These dimension-specific normal forms are not definitions for arbitrary maps
$\R^n\to\R^m$.

\section{The completed source and its five singular labels}\label{sec:atlas}
The angular formula is smooth where $c>0$ but fails to be a smooth chart at
$Q=0$. This defect is repaired by changing the source coordinate, while keeping
the full physical surface \cite{ruled}.
\begin{proposition}[Endpoint charts and complete rank set]\label{prop:atlas}
Near either endpoint let $\epsilon\in\{1,-1\}$, $r=Q\geq0$, and set
\[
 C(r)=1-\sqrt{1-r^2},\qquad
 B_\epsilon(r,u)=\bigl(C+2u(1-C),\epsilon(1-u)\sqrt{1-C^2},ur\bigr).
\]
The formula extends smoothly to $|r|<1$ near $r=0$, agrees with $S$ on the
physical overlap, and satisfies
\[
 (B_\epsilon)_r(0,u)=(0,0,u),\quad
 (B_\epsilon)_u(0,u)=(2,-\epsilon,0),\quad
 (B_\epsilon)_r\times(B_\epsilon)_u=(\epsilon u,2u,0).
\]
The rank-loss labels of the completed source are exactly
\[
 P_0=(0,1),\quad P_\pm=(\pm t_0,u_0),\quad
 E_\pm=(\pm\pi/2,0),\qquad
 c_0=\frac{3-\sqrt5}{2},\ t_0=\arccos c_0,\ u_0=\frac{1-c_0}{2}.
\]
Their smooth extensions are cross-caps. The physical restrictions at $P_0$
and $E_\pm$ are respectively a boundary half-germ and corner quarter-germs.
\end{proposition}
\begin{proof}
The radicands are positive near $r=0$. Since
$C(r)(2-C(r))=r^2$, $C(Q)=c$ when $0\leq c\leq1$; the sign of $s$ fixes
$\epsilon$. Differentiation gives the displayed endpoint vectors, so an
endpoint loses rank exactly at $u=0$. On $c>0$, the full cross-product
equations reduce either to $c=1,u=1$, or to
$u=(1-c)/2=c(2-c)/2$. The latter equation is $P(c)=c^2-3c+1=0$ and has
exactly one root in $(0,1)$. Direct substitution checks the remaining component.
Thus $P(c)$ selects the lateral points, not all surface singularities.

For a rank-one ruled germ choose the nonzero ruling vector and its kernel
direction. Because $S_{uu}=0$, the cross-cap recognition determinant reduces
to $\det(S_u,S_{ut},S_{tt})$. Its values are
\[
 \Delta(P_0)=1,\qquad \Delta(P_\pm)=-5\sqrt{\sqrt5-2}\ne0.
\]
In the endpoint chart,
$\det(B_u,B_{ur},B_{rr})(0,0)=-\epsilon\ne0$.
The smooth cross-cap criterion applies to these actual maps with their adapted
directions \cite[Corollary 4.5 and its proof]{whitneymetric}. It supplies
smooth source and target diffeomorphisms, beyond a two-jet resemblance.
The half- and quarter-germs refer to the original physical inequalities in
those charts; no coordinate straightening of their boundaries is assumed.
\end{proof}

\begin{figure}[ht]
\centering
\setlength{\unitlength}{1mm}
\begin{picture}(125,55)
\put(15,8){\framebox(95,35){}}
\put(15,43){\circle{1.7}}\put(110,43){\circle{1.7}}
\put(15,8){\circle*{1.7}}\put(110,8){\circle*{1.7}}
\put(62.5,43){\circle*{1.7}}
\put(26.8,18.8){\circle*{1.7}}\put(98.2,18.8){\circle*{1.7}}
\put(7,1){$E_-$}\put(105,1){$E_+$}\put(60,47){$P_0$}
\put(21,22){$P_-$}\put(96,22){$P_+$}
\put(54,1){$t$}\put(4,25){$u$}
\end{picture}
\caption{Filled markers are the five rank-loss source labels; open markers are
the two regular upper corners, which remain distinct source labels.
Positions are drawn from $t_0,u_0$ of
Proposition~\ref{prop:atlas}; this parameter diagram is not a surface embedding.}
\end{figure}

Put $D_{\rm reg}=D\setminus\{P_0,P_+,P_-,E_+,E_-\}$ with the completed atlas.
It is a connected orientable surface with retained regular boundary corners.
Deleting boundary points from a disk does not create an interior hole:
one may retract collar arcs into the interior. Removing the two interior
points leaves a twice-punctured disk. For example, excise two disjoint small
disks and retract the remaining disk onto two circles joined to a base point
by disjoint arcs. Hence $H_1(D_{\rm reg};\Z)=\Z^2$, and its orientation cover
is the disconnected trivial cover. Orientation comes from the source charts;
there is no need for an invalid angular normal at $Q=0$. The paired top-edge
image relation has a fixed point and supplies no free boundary gluing.

\section{The actual observation field and square-root lifting}\label{sec:field}
On the open angular chart $|t|<\pi/2$ define
\begin{equation}\label{eq:F}
 N=(1-c)(1-c-2u),\quad
 M=\frac{c^2(2-c)-u(1-c+c^2)}{Q},\quad F=N+iM.
\end{equation}
Both $N$ and $M$ are even in $t$. Their Jacobian determinant is odd, because
the $t$ derivatives reverse sign. A source area reparametrization multiplies
this field by its Jacobian. The angular field generally diverges at $Q=0$;
it is not a continuous global scalar field on the completed rectangle.
\begin{theorem}[Interior dipole and restoration]\label{thm:dipole}
The interior zeros of $F$ are exactly $P_\pm$. With the $(t,u)$ orientation,
\[
 \det DF(P_\pm)=\mp\frac{5s_0(1-c_0)^2}{Q_0},\qquad
 \operatorname{Ind}(P_\pm)=\mp1.
\]
Every continuous square-root branch transported round a sufficiently small
positive loop about either point changes sign. Two traversals restore it.
\end{theorem}
\begin{proof}
Solving \eqref{eq:F} gives the same lateral roots as above. The determinant
formula follows by substitution and $c_0^2=3c_0-1$. It is nonzero.
The actual small-circle winding of a smooth planar nondegenerate zero is
the sign of its determinant: a linear homotopy to the derivative avoids zero
on sufficiently small circles; an invertible real linear map has winding
$1$ or $-1$ according to orientation. The predecessor proof applies to
exactly \eqref{eq:F}, the same open domain and orientation \cite{ruled}.

Along a loop write $F(\gamma(s))=\exp L(s)$ using a continuous logarithm
lift. Its endpoint is $L(1)=L(0)+2\pi i n$, where $n=\pm1$. Thus
$\chi(s)=\exp(L(s)/2)$ satisfies $\chi(1)=(-1)^n\chi(0)=-\chi(0)$.
Any other continuous root divided by $\chi$ is a continuous map from a
connected interval to $\{1,-1\}$ and is constant. Repeating the loop makes
the sign positive. This is a traversal count; it does not assign a physical
clock or equate the root cover with the normalization-sheet cover.
\end{proof}

\begin{proposition}[A genuine polar fold]\label{prop:polar}
The smooth extension of $\Psi=(N,M)$ at $P_0$ is a fold. It has full-extension
local degree zero. No interior winding index is assigned on its physical
half-neighbourhood.
\end{proposition}
\begin{proof}
Use source coordinates
\[
 a=-M(t,u),\qquad b=2\sin(t/2)\sqrt{u-(1-\cos t)/2},
\]
and target coordinates $(N,M)\mapsto(-M,-N)$. Near $(0,1)$ the source
radicand is positive, and its derivative matrix is
$\left(\begin{smallmatrix}0&1\\1&0\end{smallmatrix}\right)$, which is
invertible. The inverse-function theorem gives two-sided smooth charts.
The exact identity $-N=b^2$ puts the map into $(a,b^2)$; thus no conclusion
is inferred from a Taylor jet alone. Reflection of $b$ pairs regular
preimages with opposite signs, or, equivalently, the model circle lies in a
half-plane and has zero winding. The target/source orientation changes do
not change zero degree. In historical v12 the displayed remainder for $M$
omitted $M(t,1)=-t^4/2+O(t^6)$; this term does not invalidate these charts.
\end{proof}

\section{Residual volume and an explicit complex completion}\label{sec:residual}
For the standard models take source Lebesgue measure and residuals
$K_f=x^2+y^4$ and $K_\times=x^2(1+y^2)+y^4$.
\begin{theorem}[Sublevel law with quantitative comparison]
Writing $B=B(1/4,3/2)$, one has, for $\varepsilon>0$,
\[
 V_f(\varepsilon)=B\varepsilon^{3/4},\qquad
 \frac{B\varepsilon^{3/4}}{\sqrt{1+\sqrt\varepsilon}}
 \leq V_\times(\varepsilon)\leq B\varepsilon^{3/4}.
\]
The common exponent is $3/4$. A smooth positive density changes the leading
constant by its value at the origin. Smooth invertible source and target
changes preserve the exponent, but not generally the displayed coefficient.
\end{theorem}
\begin{proof}
Slicing in $y$ gives
$4\int_0^{\varepsilon^{1/4}}\sqrt{\varepsilon-y^4}\dd y$ for the fold.
Set $y=\varepsilon^{1/4}v$ and then $z=v^4$ to obtain $B\varepsilon^{3/4}$.
For the cross-cap the integrand has the additional factor $(1+y^2)^{-1/2}$;
on this interval it lies between $(1+\sqrt\varepsilon)^{-1/2}$ and $1$.
After anisotropic scaling the density converges uniformly to its value at
the origin. A target diffeomorphism fixing zero compares squared norms by
positive constants; a source diffeomorphism has a bounded positive Jacobian.
These comparisons retain the power exponent. They do not make the standard
Euclidean coefficient an invariant of arbitrary coordinates or metrics.
\end{proof}
The corresponding geometric cost satisfies
$-\log V(\varepsilon)=\tfrac34\log(\varepsilon^{-1})-\log B+o(1)$ in the
specified standard measure. Separately, $|P'(c_0)|=\sqrt5$ for the monic
$P(c)=c^2-3c+1$; multiplying $P$ by a nonzero constant changes that derivative.

\begin{theorem}[Chosen $A_3$ completion and its cohomology]
Extend the coefficients of $K_\times$ to $\C$. Its holomorphic germ is
right-equivalent to $X^2+Y^4$, with Milnor and Tjurina numbers three.
The affine fiber $U=\{X^2+Y^4=1\}$ compactifies to a smooth genus-one
curve $E$ with two points removed. Its first cohomology has dimension three
and monodromy spectrum $\{-i,1,i\}$. Its mixed Hodge weight-graded
dimensions in weights one and two are respectively two and one.
\end{theorem}
\begin{proof}
On a disk $|y|<1$ choose the holomorphic square root of $1+y^2$ taking value
one at zero. The coordinates $X=x\sqrt{1+y^2},Y=y$ have invertible derivative
and give the equality exactly. The Jacobian algebra is
$\C\{Y\}/(Y^3)$; weighted Euler's identity puts $X^2+Y^4$ in its Jacobian
ideal, giving the same Tjurina dimension.

The weighted projective compactification is
$X^2+Y^4=Z^4$ in $\mathbb P(2,1,1)$. It avoids the ambient weighted
singular point. The map to $[Y:Z]\in\mathbb P^1$ is a double cover
branched at four simple roots. Riemann--Hurwitz gives genus one. At infinity
$W=X/Y^2$ takes the two distinct values $\pm i$, so both ends are smooth
and unramified. The first Betti number is $2g+2-1=3$. Weighted homogeneity
identifies this fiber with the local Milnor fiber \cite{milnor,dimca}.
The monodromy $h(X,Y)=(-X,iY)$ pulls $dY/X$ back to $-i\,dY/X$ and acts
by $i$ on its conjugate. Near infinity $(s,W)=(1/Y,X/Y^2)$ transforms as
$(-is,W)$, so the boundary class is fixed. Finally the Gysin sequence of
mixed Hodge structures is
\[
 0\longrightarrow H^1(E;\Q)\longrightarrow H^1(U;\Q)
 \longrightarrow\widetilde H^0(\{\infty_+,\infty_-\};\Q)(-1)
 \longrightarrow0.
\]
The first term is pure of weight one with types $(1,0),(0,1)$; the last has
type $(1,1)$ and weight two \cite{deligne,peters}. These are external
cohomological theorems with their required smooth algebraic pair explicitly
constructed. No Hodge structure is inferred from the real surface or a
Mellin parameter.
\end{proof}

\section{Three different projected curves}\label{sec:links}
\begin{proposition}[Real link and embedded lift]
For every $r>0$, $L_r=\{K_\times=r^2\}$ is a smooth circle. Its cross-cap
image has one transverse double point and is homeomorphic to $S^1\vee S^1$.
The lift $(x,xy,y^2,y)$ is a smooth embedding of the source plane.
\end{proposition}
\begin{proof}
The gradient of $K_\times$ vanishes only at zero. The compact level set is
the two graphs $x=\pm\sqrt{(r^2-y^4)/(1+y^2)}$, joined at
$y=\pm\sqrt r$, and is connected. A connected compact smooth one-manifold
is a circle. Equality of cross-cap images forces $x=x'$ and $y'=\pm y$;
distinct preimages on the link occur only at $(0,\pm\sqrt r)$.
The tangent images there are $(1,\pm\sqrt r,0)$, which are transverse.
Elsewhere the source map has rank two, so the link restriction is immersed.
The compact-domain quotient identifies exactly that pair and gives the
figure-eight graph. The first and fourth lift coordinates supply a continuous
smooth inverse, and their differential has rank two everywhere. This proves
an embedding, rather than injectivity alone.
\end{proof}

\begin{proposition}[Weighted orbit and its amplitude domain]
The primitive positive weights of $X^2+Y^4$ are $(2,1)$.
For $a,b>0$, $a^2+b^4=1$, the orbit
$\Gamma(s)=(ae^{2is},be^{is})$, $s\in\R/2\pi\Z$, is embedded.
Its rectangular shadow $(b\sin s,a\sin2s)$ is a transverse figure-eight;
adding $\Re Y$ reconstructs the orbit. The frequency ratio is $2:1$.
The quadratic phase action is $2\pi(4a^2+b^2)$. Its infimum on
nondegenerate orbits is $2\pi$, attained only in their degenerate closure;
its unique interior critical point $b^2=1/8$ is a strict maximum.
\end{proposition}
\begin{proof}
The equation $2w_X=4w_Y$ and coprimality give $(2,1)$.
The $Y$ coordinate determines the phase, so the lifted circle is embedded.
The shadow has only the distinct pair $s=0,\pi$ over its double point;
the tangent determinant is $4ab\ne0$. Eliminating phase gives
$v^2=4a^2(u^2/b^2)(1-u^2/b^2)$, and adding $b\cos s$ fixes phase.
Differentiation gives $X''+4X=0$, $Y''+Y=0$.
With $z=b^2\in(0,1)$ the action divided by $2\pi$ is
$4+z-4z^2=1+(1-z)(4z+3)>1$. Its derivative vanishes only at $z=1/8$
and its second derivative is $-8$. The value $1$ occurs at $z=1$ in the
closed family. Replacing $s$ by $\omega_0 t$ changes the clock for every
$\omega_0>0$. The orbit leaves the fixed fiber $f=1$ at general phases:
$f(\Gamma(s))=e^{4is}$. It is an orbit in the ambient completion.
\end{proof}

\begin{proposition}[Actual principal-curvature cusp]\label{prop:cusp}
For the actual lateral germs, the Euclidean Whitney invariants obey
$b_\pm^2+a_\pm c_\pm^2=0$ with $c_\pm>0$.
Their principal-curvature exceptional traces have an ordinary cusp at the
double return, rather than a transverse node.
\end{proposition}
\begin{proof}
Use the actual adapted pair
$\xi=\partial_u/\sqrt3$, $\eta=-\partial_t+\lambda_\pm\partial_u$,
$\lambda_\pm=-c_0s_\pm/[2(1-c_0)]$. Let $\Pi$ project onto the normal
plane of $S_u$, and put
$A=\Pi(S_{tt}-2\lambda_\pm S_{tu})$, $B=\Pi(-S_{tu}/\sqrt3)$.
The independent derivative computation yields
\[
 \begin{aligned}
 N=|A|^2&=\frac{-335+162\sqrt5}{12},\\
 M^2=\langle A,B\rangle^2&=\frac{-520+249\sqrt5}{27},\\
 R=|B|^2-M^2/N&=\frac{20(140+11\sqrt5)}{11397}>0.
 \end{aligned}
\]
Oriented Gram--Schmidt gives
$a=-M^2/(2N^{3/2})$, $b_\pm=\pm\sqrt{M^2R}/N$,
$c=\sqrt{2R}/N^{1/4}$. Hence the discriminant vanishes exactly.
The external Euclidean normal-form and pedal identities
\cite[Proposition 2.1 and Eq.~(2.8)]{principal} apply to these nondegenerate
cross-caps. In their notation the pedal trace is
\[
 p(\tau)=\frac{2(c\tau^2+2b\tau-ac)}{c^2+4\tau^2}(-c,2\tau).
\]
Set $\tau_*=-b/c$. Its numerator is $c(\tau-\tau_*)^2$.
Writing $p(\tau_*+h)=h^2v(h)$, one has $v(0)\ne0$ and
\[
 \det(p''(\tau_*),p'''(\tau_*))
 =-\frac{96c^3}{(c^2+4b^2/c^2)^2}\ne0.
\]
This proves the ordinary cusp criterion, beyond the vanishing discriminant.
It concerns the exceptional trace specified by the pedal theorem; it is not
a global classification of the whole principal-curvature surface.
\end{proof}

\begin{table}[ht]\centering
\begin{tabular}{@{}lll@{}}\toprule
Construction & Source & Projected feature\\\midrule
Real cross-cap link & $K_\times=r^2$ & transverse node\\
Weighted phase orbit & circle in $\C^2$ & transverse node\\
Actual exceptional trace & adapted direction $\tau$ & ordinary cusp\\\bottomrule
\end{tabular}
\caption{Distinct maps produce distinct curves. No amplitude, clock, or
scattering rule transfers between these constructions.}
\end{table}

\section{Mellin integration, measures and operator domains}\label{sec:mellin}
\begin{theorem}[Formal Green identity and measure dependence]
Let $T_\sigma f=-i(xf'+\sigma f)$, $\sigma\in\R$. For complex $C^1$
functions on $[a,b]\subset(0,\infty)$,
\[
 \int_a^b\bigl(\overline{T_\sigma f}\,g-\overline f\,T_{1-\sigma}g\bigr)\dd x
 =i[x\overline f g]_a^b.
\]
Thus on $C_c^\infty((1,\infty))\subset L^2(dx)$ the formal adjoint
expression is $T_{1-\sigma}$ and symmetry holds exactly at $\sigma=1/2$.
For measure $x^{2\beta-1}\dd x$, the formal adjoint parameter is
$2\beta-\sigma$, so the symmetry parameter is $\sigma=\beta$.
\end{theorem}
\begin{proof}
The first integrand is $i(x\overline f g)'$; integrating gives the endpoint
term. Compact support permits an interval with zero endpoint terms and
then the actual improper integral equality. The compact core is dense by
local truncation and smooth approximation in $L^2$. On a weighted space,
$ (x x^{2\beta-1})'/x^{2\beta-1}=2\beta$ replaces the constant one in
the same integration by parts. Necessity of the formal coefficient equality
can be tested on a nonzero compactly supported $f=g$. The weighted inner
product in the historical introduction therefore cannot be silently used
for the Lebesgue threshold: if $\beta=\sigma$, the weighted expression is
formally symmetric for every $\sigma$.
\end{proof}

\begin{proposition}[Actual half-line closures and adjoints]
For $\sigma=1/2$ define
$(Uf)(s)=e^{s/2}f(e^s)$ from $L^2([1,\infty),dx)$ to $L^2([0,\infty),ds)$.
It is unitary and takes the compact-core operator to $-i\partial_s$.
Its closure has domain $H^1_0(0,\infty)$, its adjoint has domain
$H^1(0,\infty)$, and its deficiency indices are $(1,0)$.
On the full multiplicative line the corresponding closed domain is carried
to $H^1(\R)$; that operator is self-adjoint.
\end{proposition}
\begin{proof}
The substitution $x=e^s$ proves norm equality and supplies the inverse;
the chain rule gives the operator identity. The graph norm of $-i\partial_s$
is exactly the $H^1$ norm, so closure of the compact core is $H^1_0$.
Distributional integration by parts characterizes the adjoint domain as
$H^1$ with no endpoint constraint. The adjoint Green term is
$-i\overline{h(0)}k(0)$ under our convention; there is no term at infinity
for $H^1$ functions. The equations $-ih'=\pm ih$ give $h=e^{-s}$ and
$h=e^s$; only the first is square-integrable. Unequal indices exclude every
self-adjoint extension. On $\R$ neither solution is square-integrable,
and the Fourier transform identifies the closure with multiplication by
real frequency on its maximal domain \cite{reed}.

Alternatively the Haar-to-Lebesgue map $Wf(x)=x^{-1/2}f(x)$ is unitary
from $L^2(dx/x)$ to $L^2(dx)$ on the full group and conjugates
$-ix\partial_x$ to $T_{1/2}$. Restriction to a half-line does not preserve
the full-line self-adjointness result.
\end{proof}
Direct differentiation gives $T_\sigma x^{-\sigma-i\omega}
=-\omega x^{-\sigma-i\omega}$. Its half-line norm squared is
$\int_1^\infty x^{-2\sigma}\dd x$, finite exactly for $\sigma>1/2$.
Even then its trace at $x=1$ is one, so it is outside the minimal zero-trace
domain. At $\sigma=1/2$ it is a generalized eigenfunction, not an $L^2$
vector. Real formal eigenvalues alone imply no orthogonality.

\section{Actual finite-time Gram integrals}\label{sec:gram}
Let $\omega_1,\ldots,\omega_M$ be distinct real frequencies, $T>0$,
and $\psi_j(t)=e^{-i\omega_jt}$. Set
\[
 G_{jk}=T^{-1}\int_0^T\overline{\psi_j(t)}\psi_k(t)\dd t.
\]
This convention makes the quadratic form $c^*Gc$ the norm squared of the
synthesis $\sum_jc_j\psi_j$, divided by $T$.
\begin{theorem}[Independence and quantitative spectral enclosure]
For $M\geq2$ put $\delta=\min_{j\ne k}|\omega_j-\omega_k|>0$ and
$\varepsilon=2(M-1)/(T\delta)$. Then
\[
 G_{jj}=1,\quad
 G_{jk}=\frac{e^{iT(\omega_j-\omega_k)}-1}{iT(\omega_j-\omega_k)},\quad
 |G_{jk}|\leq\frac{2}{T|\omega_j-\omega_k|}\ (j\ne k),\quad
 \norm{G-I}_{2\to2}\leq\varepsilon.
\]
The phasors are independent and $G$ is positive definite for every $T>0$.
If $\varepsilon<1$, then
\[
 1-\varepsilon\leq\lambda_j(G)\leq1+\varepsilon,\qquad
 \kappa_2(G)\leq\frac{1+\varepsilon}{1-\varepsilon},\qquad
 M\log(1-\varepsilon)\leq\log\det G\leq0.
\]
In particular $-\log\det G\leq M\varepsilon/(1-\varepsilon)$.
For $M=1$, $G=(1)$; an empty family has empty Gram matrix and determinant
one, with no condition number or separation minimum required.
\end{theorem}
\begin{proof}
The entries are actual elementary integrals. The numerator has modulus at
most two. The matrix is Hermitian. Each row of $G-I$ has absolute sum at
most $\varepsilon$, so Gershgorin places all its eigenvalues in
$[-\varepsilon,\varepsilon]$. The spectral theorem for a finite Hermitian
matrix then gives the Euclidean operator norm, not just an arbitrary
matrix norm.

For independence, a synthesis vanishing almost everywhere on $[0,T]$
vanishes everywhere by continuity. Differentiate at any interior time
$t_*$ to orders $0,\ldots,M-1$. The coefficient vector
$(c_je^{-i\omega_jt_*})_j$ is killed by the Vandermonde matrix of the
distinct numbers $-i\omega_j$. Its determinant is nonzero, so $c=0$.
This proves positivity for all positive $T$; the separation threshold is
only a sufficient conditioning estimate. Multiplying the eigenvalue lower
bounds gives the log-determinant lower bound. Hadamard's inequality applied
to the positive Gram matrix with unit diagonal gives $\det G\leq1$.
Finally $-\log(1-\varepsilon)\leq\varepsilon/(1-\varepsilon)$ follows by
integrating $(1-s)^{-1}$ on $[0,\varepsilon]$.
\end{proof}
For distinct positive integers take $\omega_j=\log n_j$. Multiplicative
independence is unnecessary. These estimates compare a fixed finite family
as $T\to\infty$; they are not a uniform theorem for growing frequency sets.

\section{Hardy evaluation and Bohr--Haar completion}\label{sec:hardy}
Define the Dirichlet Hilbert space by its coefficient norm:
$\mathcal H_D^2=\{(a_n):\sum|a_n|^2<\infty\}$, with Dirichlet evaluation
initially defined on finite polynomials. For $\Re s=\sigma>1/2$,
Cauchy--Schwarz gives
$|\sum a_nn^{-s}|\leq\norm{a}_{\ell^2}\zeta(2\sigma)^{1/2}$.
At $\sigma\leq1/2$, use the finite coefficient vector proportional to
$(n^{-\bar s})_{n\leq N}$. Its unit-norm evaluation is
$(\sum_{n\leq N}n^{-2\sigma})^{1/2}\to\infty$.
Thus evaluation has no bounded extension at or below that boundary.
This direct proof also specifies what ``evaluation'' means where not every
infinite series converges \cite{hardy}.

Write $n=\prod_p p^{v_p(n)}$ and
$\chi_n(z)=\prod_pz_p^{v_p(n)}$ on the product torus $\T^{\mathcal P}$.
With product normalized Haar measure,
\[
 \int\overline{\chi_n}\chi_m\dd m_H=\delta_{nm},\qquad
 \lim_{T\to\infty}T^{-1}\int_0^T n^{it}m^{-it}\dd t=\delta_{nm}.
\]
The first identity follows by integrating the finite coordinate product:
unique factorization makes at least one nonzero exponent when $n\ne m$,
and its circle integral is zero. The second follows from the actual Gram
integral. A finite integer set may generate a lower-rank exponent lattice
(for example $2,4$); this relation is retained by the torus construction.
The half-density and evaluation boundary are two different analytic results,
not a common arithmetic or physical theorem.

\section{Induced metric, exact refinement and zero capacity}\label{sec:capacity}
On the corrected $D_{\rm reg}$ take $g=S^*\delta_{\R^3}$ in its completed
atlas, $\mathcal H_g=L^2(dA_g)$, and
\[
 \mathcal E(f,h)=\int g^{ab}\overline{\partial_af}\partial_bh\dd A_g.
\]
For nested finite-energy conforming spaces whose bases satisfy
\[
 \phi_i^{(n)}=\sum_\alpha J_{\alpha i}\phi_\alpha^{(n+1)},
\]
sesquilinearity gives $J^*G_{n+1}J=G_n$ and $J^*K_{n+1}J=K_n$ exactly.
These identities require exact integration; they assert no discretization
convergence or solver accuracy.

\begin{theorem}[Capacity and form removability]\label{thm:capacity}
Every one of the five rank-loss labels has local $1$-capacity zero, using
a full, half-, or quarter-neighbourhood as appropriate. Deleting $P_\pm$
does not change the closed form with fixed outer Dirichlet boundary.
\end{theorem}
\begin{proof}
For the standard cross-cap,
\[
 g=\begin{pmatrix}1+y^2&xy\\xy&x^2+4y^2\end{pmatrix},\quad
 w^2=\det g=x^2+4y^2+4y^4,\quad
 \rho^2=x^2(1+y^2)+y^4.
\]
Since $\rho$ is ambient distance, $|\nabla_g\rho|\leq1$ on the regular
locus. For $0<r<1$, the set $\rho<r$ lies in
$|x|<r,|y|<\sqrt r$ and $w\leq3\sqrt r$, so its area $A(r)\leq12r^2$.
Choose a cutoff equal to one below $\varepsilon^2$, zero above
$\varepsilon$, and $\log(\varepsilon/\rho)/\log(1/\varepsilon)$ between.
Smooth its corners on comparable annuli. Its squared $L^2$ norm is
$O(\varepsilon^2)$, and its energy is bounded by
\[
 \frac{C}{\log^2(1/\varepsilon)}
 \int_{\varepsilon^2}^{\varepsilon}r^{-2}\dd A(r)
 \leq\frac{C'}{\log(1/\varepsilon)}\longrightarrow0.
\]
The estimate follows from Stieltjes integration by parts:
the endpoint contribution is bounded and
$2\int r^{-3}A(r)\dd r=O(\log(1/\varepsilon))$.
These cutoffs are smooth source functions because they are constant near
the singularity. Restricting them to half- or quarter-domains decreases
their norms and energies.

The actual cross-cap equivalences in Proposition~\ref{prop:atlas} compare
the induced quadratic forms by positive constants: the target derivative
and its inverse are bounded. In dimension two this also compares area
densities and Dirichlet energies. Thus capacity transfers to all five actual
germs; no metric isometry is required. Multiplying a compactly supported
smooth test function by $1-\chi_\varepsilon$ near $P_\pm$ approximates it
in $\mathcal E+L^2$ norm, by the product rule and the displayed bounds.
The reverse closure inclusion is immediate. This is equality of form
closures, not of operator graph closures.
\end{proof}

\begin{theorem}[Actual compactness and positive Dirichlet gap]\label{thm:compact}
The fixed-outer-Dirichlet form domain embeds compactly into $\mathcal H_g$.
Its Friedrichs operator $H_F$ has compact resolvent and satisfies
$H_F\geq\lambda_D I$ for some $\lambda_D>0$.
\end{theorem}
\begin{proof}
Use the original completed-source coordinates at each of the five points,
translated to the origin, and write $R=(x^2+y^2)^{1/2}$.
The standard cross-cap has area density comparable to $R$.
The actual source and target diffeomorphisms from
Proposition~\ref{prop:atlas}, with bounded derivatives and inverses, therefore
give $\dd A_g\asymp R\dd x\dd y$ in these coordinates.
Also $g\leq C I$ since the actual extended map has bounded derivative.
Consequently the energy dominates $c\int R|\nabla f|^2\dd x\dd y$.
The physical coordinate domains are a disk, half-disk or quarter-disk;
their radial angular intervals are fixed. No metric cone equivalence is used.

Let $\chi$ be a chart cutoff, equal to one near zero and vanishing near its
outer circular edge. For a smooth core function set $h=\chi f$.
On each radial ray, integration of $(R|h|^2)'$ and Cauchy--Schwarz give
\[
 \int_0^{R_0}|h|^2\dd R
 =-2\Re\int_0^{R_0}R\overline h\,\partial_Rh\dd R
 \leq2\Bigl(\int|h|^2\dd R\Bigr)^{1/2}
          \Bigl(\int R^2|\partial_Rh|^2\dd R\Bigr)^{1/2}.
\]
The endpoint terms vanish; the initial core avoids the singularity.
Squaring and integrating the fixed angular interval yields
\[
 \int\frac{|\chi f|^2}{R^2}\dd A_g
 \leq C\int R|\nabla(\chi f)|^2\dd x\dd y
 \leq C'(\mathcal E(f,f)+\norm f^2).
\]
The cutoff derivative lies on a regular compact annulus, where the last
bound follows from the product rule. The inequality extends by form closure.
In particular, the $L^2$ mass below $R<\varepsilon$ is bounded by
$C\varepsilon^2$ times the form norm squared, uniformly at all five points.

Away from these disks the actual metric is smooth and uniformly elliptic
on a compact source with ordinary Lipschitz boundary and retained regular
corners. The usual Rellich theorem \cite{brezis} gives compactness there.
A diagonal subsequence followed by the uniform small-disk mass bound gives
compactness on the whole source, including both endpoint corners.

If no positive gap existed, take $\norm{f_n}=1$ and
$\mathcal E(f_n,f_n)\to0$. Compactness gives a strong $L^2$ limit of norm
one. On every regular chart its weak derivative is zero, hence it is
constant on the connected interior. On a regular outer-boundary collar,
the zero Dirichlet trace passes to the weak $H^1$ limit and makes this
constant zero, a contradiction. This proves the gap. The form representation
theorem and compact form embedding give the claimed compact resolvent.
\end{proof}

For the standard metric an additional exact identity is
\[
 1-|\nabla_g\rho|^2
 =\frac{x^2y^4}{(x^2+4y^2+4y^4)(x^2(1+y^2)+y^4)}\leq\rho/4.
\]
The front coordinates
$x=r\cos\theta/\sqrt{1+r\sin\theta}$,
$y=\epsilon\sqrt{r\sin\theta}$, $0<\theta<\pi$, have $\rho=r$.
Each front has limiting length $\pi r$. Indeed its normalized speed
converges to one away from the endpoints and is bounded by
$C(1+\sqrt{r/\sin\theta})$, with an integrable dominating function.
The two fronts join on the source to one circle, so $L(r)/r\to2\pi$.
Wirtinger's inequality on that circle gives
\[
 \int_{\rho<r_0}\frac{|u-\Pi_\rho u|^2}{\rho^2}\dd A_g
 \leq C\mathcal E(u,u)
\]
for the standard model: coarea uses $|\nabla\rho|$ bounded above and below,
and the tangential part of the energy controls the link derivative.
In seam coordinates $x=\xi,y=\xi z$ the determinant is
$\xi^4(1+4z^2+4\xi^2z^4)$, and the density-weighted inverse metric has
Schur complement $\sqrt{1+4z^2+4\xi^2z^4}/[\xi^2(1+4z^2)]$.
On fixed strips $|z|\leq K$ it bounds the weighted one-dimensional
Poincare estimate. The transverse zero solutions, a constant and a
primitive of $\sqrt{1+4z^2}$, are not localized in
$L^2(\R,\sqrt{1+4z^2}\dd z)$.
These calculations establish model inequalities. They alone do not exclude
all defect modes of an actual patched unbounded operator.

\section{The operator domain and its boundary obstruction}\label{sec:green}
Let $H_F$ be the self-adjoint operator associated with the closed fixed-outer-
Dirichlet form. This definition uses form closure and makes sense before a
point-trace theorem is available. Historical v12 instead set the punctured
minimum equal to the graph closure of $-\Delta_g$ on
$C_c^\infty(D^\circ\setminus\{P_+,P_-\})$. Denote that closure by $H_0$.
\begin{theorem}[Outer-boundary deficiency obstruction]\label{thm:outer}
For that definition of $H_0$, both deficiency spaces are infinite-dimensional.
In particular, its claimed indices $(2,2)$ and exhaustive $U(2)$ extension
family are false. Fixing the outer boundary in a sentence does not alter
the adjoint of the specified compact-core operator.
\end{theorem}
\begin{proof}
Choose a regular open outer-boundary arc away from all five singular labels
and corners. It carries a smooth uniformly elliptic metric and the usual
Sobolev trace. For any smooth boundary value $h$ supported inside this arc,
choose a smooth collar extension $w_h$, supported away from the singular
points. For $z\in\C\setminus\R$ put
\[
 v_h=w_h-(H_F-z)^{-1}(-\Delta_g-z)w_h.
\]
This is an $L^2$ function. Distributionally on the interior it satisfies
$(-\Delta_g-z)v_h=0$, so testing against the punctured compact core puts
$v_h$ in $\ker(H_0^*-z)$. The resolvent term belongs to the fixed Dirichlet
form domain and has zero trace on the regular arc. Hence $v_h$ has trace
$h$. Linearly independent choices of $h$ give linearly independent $v_h$.
The argument works for $z=i$ and $z=-i$ and requires no claim about the
cross-cap Green asymptotics. The Hilbert space is separable, so these
deficiency dimensions are countably infinite.

An elementary interval witness illustrates the graph-closure distinction:
$v(x)=x(1-x)$ has Dirichlet values zero but
$\langle-v'',1\rangle-\langle v,-1''\rangle=2$. The same pairing is zero
for every interior compact test function and remains zero under graph
limits. Thus even the ordinary Dirichlet domain is larger than that compact
graph closure.
\end{proof}

\begin{theorem}[Conditional point restriction with fixed outer boundary]\label{thm:conditional}
Assume that the actual $H_F$ domain admits a graph-continuous surjective
trace $\tau:\Dom H_F\to\C^2$ at the interior singularities, and that
$\ker\tau$ is dense in $\mathcal H_g$. Then
$A=H_F|_{\ker\tau}$ is closed, densely defined and symmetric, with
deficiency indices $(2,2)$. Its self-adjoint extensions are parametrized by
$U(2)$. These hypotheses do not assert the actual existence of the traces.
\end{theorem}
\begin{proof}
The kernel is closed in the graph norm. For nonreal $z$, the bounded
surjection $\tau(H_F-z)^{-1}:\mathcal H_g\to\C^2$ has kernel
$\operatorname{ran}(A-z)$. The adjoint of this map supplies two independent
defect vectors, because
$\ker(A^*-\bar z)=\operatorname{ran}(A-z)^\perp$. Thus both dimensions
are two. Von Neumann's unitary classification now applies to this corrected
operator. The regular Dirichlet point-interaction construction
\cite[Proposition 2.1 and Theorem 2.2]{point} uses exactly this restriction,
but its interior elliptic $H^2$ point estimate is not a theorem for our
degenerate Whitney metric. The proof here is an abstract conditional theorem.
\end{proof}

The standard metric has the exact residual
\[
 \Delta_g\log\rho=
 \frac{2x^2y^2(3x^2y^4+5x^2y^2+x^2+7y^6+5y^4)}
 {(x^2+4y^2+4y^4)^2(x^2(1+y^2)+y^4)^2}.
\]
To check integrability directly, set $R=(x^2+y^2)^{1/2}<1$.
The numerator is at most $C x^2y^2(x^2+y^4)$, whereas the denominator
is at least $(x^2+y^2)^2(x^2+y^4)^2$. Thus the residual is bounded by
$C/R^2$ and $\dd A_g\leq C R\dd x\dd y$.
Its $p$th power is bounded in integral by
$C\int_0^{R_0}R^{2-2p}\dd R$, finite for $1<p<3/2$.
On $\rho=r$, the preceding radial defect bound and link length limit give
$r^{-1}\int_{\rho=r}|\nabla_g\rho|\dd\ell\to2\pi$.
This standard flux and residual motivate a logarithmic Green
candidate. It does not prove that applying an actual resolvent to the
remainder gives a continuous correction, a graph-bounded point trace,
or an exhaustive maximal-domain expansion. Those are different estimates.
The required completion is a two-sided parametrix in specified front/seam
weighted spaces, with an overlap estimate controlling the localized graph
norm and proving the actual resolvent maps that $L^p$ remainder into
continuous functions. A global parametrix must also cover both endpoint
corner charts, although compactness is already established in
Theorem~\ref{thm:compact}. The conifold assumptions in \cite{pacini} do not hold
automatically in the degenerating seam coordinates.

\begin{proposition}[Reference-scale invariant boundary planes]
On $\C^2_a\oplus\C^2_b$ with form
$\omega((a,b),(c,d))=2\pi(\langle a,d\rangle-\langle b,c\rangle)$,
the only complex Lagrangian plane invariant under every
$R_t(a,b)=(a,b+ta)$ is $\{a=0\}$.
\end{proposition}
\begin{proof}
For $(a,b)$ in an invariant plane, subtract its $R_1$ image to see that
$(0,a)$ belongs to the plane. Isotropy gives $2\pi|a|^2=0$, so $a=0$.
Dimension equality gives the stated plane; direct substitution checks
invariance. This is a finite-plane theorem. Identifying it with the
Friedrichs extension requires the actual Green coefficients and form/domain
comparison, beyond Theorem~\ref{thm:conditional}.
\end{proof}
Zero puncture capacity proves the fixed-boundary form removability of
Theorem~\ref{thm:capacity}. A claim of uniqueness among all submarkovian
extensions additionally needs a theorem classifying those extensions by the
relevant Dirichlet forms with the same outer boundary. The old minimum admits
other outer-boundary conditions, so that unrestricted uniqueness claim also
fails. Concretely, close the same energy on smooth completed-source functions
allowed to have nonzero outer trace. Such functions have finite local energy:
bounded derivatives and the cross-cap metric give an integrable $O(1/R)$
energy density in source coordinates. The form is closable, as can be checked
first on regular compact subdomains and then in the global weighted gradient
space. Smooth contractions and the chain rule give the Markov property after
closure. The associated Neumann form operator extends the interior compact
operator and contains the constant function in its zero eigenspace. In
contrast, the positive Dirichlet gap excludes that eigenfunction for $H_F$.
This supplies two distinct submarkovian extensions of the old $H_0$.
Their existence is not altered by the finite-plane calculation.

\begin{theorem}[Actual Krein and buckling results for the compact-core minimum]
The actual closed compact-core operator $H_0$ is densely defined, symmetric
and strictly positive. Its Friedrichs extension is $H_F$.
Its Krein extension satisfies
\[
 \Dom(H_0)_K=\Dom H_0\dotplus\ker H_0^*,\qquad
 \ker(H_0)_K=\ker H_0^*.
\]
The kernel is infinite-dimensional. The reduced positive spectrum is
discrete with finite multiplicities, and its eigenvalues are the weak
buckling eigenvalues on $\Dom H_0$:
\[
 \langle H_0v,H_0\varphi\rangle
 =\lambda\langle v,H_0\varphi\rangle\qquad(\varphi\in\Dom H_0).
\]
\end{theorem}
\begin{proof}
Interior compact smooth functions are dense in $\mathcal H_g$ by truncation
and smooth approximation on the regular locus. Integration by parts makes
the initial operator symmetric, and its graph closure is contained in
$H_F$. The positive gap in Theorem~\ref{thm:compact} passes to this closure.
The capacity theorem identifies the two Friedrichs forms, hence its
Friedrichs extension is $H_F$.
For the kernel repeat the collar construction in
Theorem~\ref{thm:outer} at $z=0$; the bounded inverse $H_F^{-1}$ now
exists. Independent boundary values inject into $\ker H_0^*$.
The abstract Krein/buckling theorems apply to this actual closed densely
defined strictly positive operator \cite[Hypothesis 2.2, Theorem 2.4 and
Theorem 3.4]{buckling}. Compact resolvent of $H_F$ gives discreteness of
the reduced positive Krein spectrum. This retains the valid global spectral
conclusion while correcting the false finite-channel dimension.
\end{proof}

Under the point-trace hypotheses of Theorem~\ref{thm:conditional}, the
corrected $A$ is also strictly positive by restriction, and the abstract
Krein theorem gives
\[
 \Dom A_K=\Dom A\dotplus\ker A^*,\qquad\ker A_K=\ker A^*.
\]
If actual logarithmic boundary coordinates and zero-energy Green vectors
are established, this domain becomes $b=R_0a$, with $R_0$ Hermitian and
reference transformation $R_0\mapsto R_0+tI$. Its reduced nonzero
spectrum obeys the weak buckling equation
$\langle Av,A\varphi\rangle=\lambda\langle v,A\varphi\rangle$ on
$\Dom A$ \cite{krein,buckling}. Its Friedrichs form must be compared with
the fixed Dirichlet form before transferring that compactness conclusion to
$A$. The actual point traces, finite trace rank and logarithmic coordinates
remain unproved. The global compactness result above concerns $H_F$ and the
explicitly defined $H_0$; it does not supply these missing Green coefficients.

\section{Normalization defect and the support obstruction}\label{sec:sheaf}
Let $X$ be a sufficiently small standard cross-cap image, and
$\nu:\widetilde X\to X$ its proper topological normalization, obtained from
the source map. It is one-to-one off the open double ray $L$ and two-to-one
on $L$, with one preimage at the pinch. In the ordinary sheaf category,
\[
 \mathcal Q_\nu=\nu_*\Q_{\widetilde X}/\Q_X
\simeq j_!\mathcal L_{\rm sgn},\qquad j:L\hookrightarrow X.
\]
Here $j$ is locally closed; extension by zero means the open-ray extension
inside its closed support, followed by the closed direct image.
Indeed the diagonal stalk map is an isomorphism at one-preimage points and
has one-dimensional anti-invariant quotient at two-preimage points.
Properness gives these direct-image stalks by shrinking neighbourhoods.
The quotient pinch stalk is zero. The restriction is the sheet-sign line;
the extension-by-zero description follows by the adjunction map and its
stalk isomorphisms \cite{stacks}.

Let $i$ include the pinch. Every ordinary morphism
$j_!\mathcal L_{\rm sgn}\to i_*\C$ is zero, since every stalk map is zero.
For the local ray the nearby group $H^0(i^*Rj_*\mathcal L_{\rm sgn})$
and degree-one costalk $H^1(i^!j_!\mathcal L_{\rm sgn})$ are one-dimensional.
One may compute the latter from the compactly supported cohomology of
$(0,\varepsilon)$, or the local recollement triangle. Sheet exchange acts
by $-1$. A linear equivariant map from that line to a sheet-even line
must satisfy $T(-v)=T(v)$ and therefore vanishes in characteristic zero.
This sheaf theorem is independent of the existence of an actual scalar Green
line. The standard logarithm is sheet-even because
$\rho(0,y)=\rho(0,-y)$; exporting this interpretation to an actual exhaustive
Green space still depends on the analytic construction above.

An underlying compact or condensed set does not fix a structure sheaf:
even a one-point space supports both $\Q$ and $\Q[\epsilon]/(\epsilon^2)$.
The chosen algebraic fiber instead has its explicitly supplied coordinate
ring $\Q[X,Y]/(X^2+Y^4-1)$. Neither its structure sheaf nor its mixed Hodge
data is recovered from a bare topological normalization defect.

\section{Phase lines, cone models and transmission hypotheses}\label{sec:spin}
The squaring cover pulled back by $F/|F|$ supplies a flat Hermitian sign
line on each punctured interior source disk. Theorem~\ref{thm:dipole}
gives its holonomy $-1$, with no continuous holonomy parameter.
On an exact round cone the scalar twisted angular eigenvalues are
$n+1/2$, and the radial indicial powers are $r^{\pm|n+1/2|}$.
Since $r^\lambda\in L^2(r\dd r)$ exactly for $\lambda>-1$, the two
weights $n=0,-1$ are limit-circle channels. This scalar model has two
lowest channels, not a preferred one-dimensional normalization defect.

For the disk-bounding spin structure the untwisted angular weights are also
$\Z+1/2$. Twisting by the sign line makes them $\Z$.
The chiral exact-cone operator has the form
$c(dr)(\partial_r+1/(2r)+A_\partial/r)$. On its angular zero eigenspace,
$r^{-1/2}f(r)$ gives a unitary map from $L^2(dr)$, reducing the operator
to $c(dr)\partial_r$. Locally at the tip, $H^1$ has an endpoint trace and
the minimal closure has zero trace. Thus that exact block contributes one
complex endpoint coefficient per chirality. The full block contributes two.
This is a local tip quotient after excluding the regular outer-boundary
trace, not the unrestricted maximal/minimal quotient on a bounded punctured
disk. The latter contains infinitely many regular boundary traces.

The intrinsic periodic Dirac spectrum on any metric link circle of length
$L(r)$ is $2\pi\Z/L(r)$ once its total flat holonomy is trivial. This
link result does not construct a parametrix for the surface Dirac operator.
The model inverse for $\partial_r+k/r$ is
\[
 Q_kf(r)=\begin{cases}
 r^{-k}\int_0^r s^kf(s)\dd s,&k>0,\\
 -r^{-k}\int_r^{r_0}s^kf(s)\dd s,&k<0,\\
 \int_0^rf(s)\dd s,&k=0.
 \end{cases}
\]
Differentiation checks $(\partial_r+k/r)Q_kf=f$. To patch such inverses,
one must prove boundedness in chosen graph spaces and control the overlap
remainders. The claim that all cutoff commutators gain positive weight is
false without further construction: under $r=\varepsilon R$,
$\varepsilon\partial_r\chi(r/\varepsilon)=\chi'(R)$. Composing that
commutator with a first-order model inverse leaves weight zero, not a
factor tending to zero. Shrinking alone therefore does not justify the
claimed Neumann series. Cone-operator theory \cite{cone,dirac} cannot be
imported until the actual Whitney seam and domain hypotheses are verified.
The actual chiral/full trace dimensions in v12 are consequently unproved.

\begin{proposition}[First-order phase/deck comparison]
The pullback $h_F(v,w)=\Re(dF(v)\overline{dF(w)})$ is positive definite at
$P_\pm$. Its restriction normalizes the kernel of $dS$.
For $k_\pm=\partial_t+c_0s_\pm\partial_u/[2(1-c_0)]$,
\[
 dS(k_\pm)=0,\qquad |dF(k_\pm)|^2=655/4-72\sqrt5>0.
\]
The two normalization branches have opposite first-order $F$ phases.
Their four limiting square roots form a nonsplit $\Z_4$ torsor over the
two branches. A chosen orthogonal source reflection lifts to Pin and
exchanges the two chiral Clifford lines.
\end{proposition}
\begin{proof}
Invertibility of $dF$ proves positive definiteness. Actual differentiation
gives $dF(k_\pm)=s_\pm(1-2c_0)
-i s_\pm(15c_0-4)/[2Q_0(1-c_0)]$; reduction by $c_0^2=3c_0-1$
gives the squared norm. The two double branches approach in directions
$\pm k$, so their normalized phases tend to $a,-a$. If $\alpha^2=a$,
their roots are $\{\alpha,i\alpha,-\alpha,-i\alpha\}$. Multiplication by
$i$ squares to the phase-deck sign, and the extension
$\Z_2\to\Z_4\to\Z_2$ has no group section. Finally Clifford
multiplication by a unit reflecting vector anticommutes with chirality,
and its two Pin lifts differ by a central sign \cite{pin}.
\end{proof}
The reflection requires a chosen positive metric and either a Pin$^+$ or
Pin$^-$ convention. It is not an automatic isometry of the actual induced
metric. Right--left cross-cap equivalence only compares metrics: for example
the target shear $(X,Y,Z)\mapsto(X,Y,Z+Y)$ preserves the cross-cap germ
but destroys the standard $y\mapsto-y$ source metric symmetry.
Likewise a Pin map of the spin fibers does not supply an identification of
the separate phase-line fibers. Their first-order quarter-turn is precisely
the distinction that the four-phase torsor records.

\begin{proposition}[Conditional full-spin cut action]
On regular paired cut banks assume specified trace spaces, Dirac Green's
formula, and a unitary full-spin bundle map $U$ including a phase-line
identification, such that $U^*c(\nu_2)U=-c(\nu_1)$.
Then the graph $(v,\lambda Uv)$ is isotropic for the Green form exactly
when $|\lambda|=1$, and is fiberwise maximal isotropic.
The real Green transgression coefficient is uniquely $1/2$.
If a fixed Pin reduction is additionally required, scalar multiples of its
chosen lift remain Pin lifts exactly for $\lambda=\pm1$.
\end{proposition}
\begin{proof}
The Green pairing restricted to the graph is
$(1-|\lambda|^2)\int\langle c(\nu_1)v,w\rangle\dd s$.
The boundary Clifford map is invertible, so isotropy is equivalent to the
coefficient being zero. The graph has half the doubled fiber dimension,
giving fiberwise maximality. For
$x=\sum_j\langle\psi_j,D_j\psi_j\rangle$, the self-pairing Green term is
$\bar x-x$. Hence $x+a(\bar x-x)$ is real for arbitrary $x$ exactly when
$a=1/2$. Its variation has no boundary term on a common isotropic trace
space. This does not single out this graph among all isotropic conditions.
The kernel of the chosen Pin double cover is $\{1,-1\}$, giving the last
claim once that reduction has been supplied \cite{dirac,pin}.
\end{proof}
The conditional trace theorem neither identifies the phase and normalization
lines nor fixes their relative unitary trivialization. Without that additional
input a continuous phase freedom remains. It also leaves an overall action
multiplier and a clock unspecified. End behaviour of the singular cut must
be included before turning the regular-bank calculation into a global
self-adjoint Dirac realization.

For the self-adjoint circle operator $-i\partial_\theta$ with
$H^1$ boundary condition $f(2\pi)=-f(0)$, the eigenvalues are
$n+1/2$. The involution $n\mapsto-n-1$ pairs each with its negative.
Thus its eta series vanishes absolutely for $\Re s>1$ and its meromorphic
continuation has value zero at $s=0$ \cite{aps}. This elementary model
produces no net spectral chirality.

\section{Structural limits and the remaining research problem}\label{sec:limits}
Several negative conclusions need little machinery but prevent unjustified
identifications. A residual does not fix communication capacity: on identical
binary labels a constant stochastic channel has capacity zero and a noiseless
one has capacity $\log2$. The kernel and resource constraints are absent from
the residual. A proposed sphere-volume to disk-area ratio with
$r=\phi^{-2}R$ is $2\pi\phi^4R$ and scales as a length; dividing by $R$
defines a dimensionless number but no mass operator.
For the canonical Kahler form on $\C^2$, the power coordinates depend on
disjoint canonical pairs and satisfy $\{|z_+|^2,|z_-|^2\}=0$.
The full $SU(2)$ moment map is a different observable algebra.
A nonempty proper subset of $SU(2)$ cannot be invariant under its full
transitive left action; this conclusion is conditional on an actual proper
source lift, which \eqref{eq:S} alone has not specified.

The chosen complex residual supplies the smooth algebraic pair $(E,D)$
needed for its mixed Hodge calculation. Real signal topology and the number
$\sigma=1/2$ do not specify such a pair, a finite-field model, Frobenius,
or an $\ell$-adic sheaf. A single curve is not a Weil cohomology theory.
For a separately specified realization extended to open varieties with
localization and Tate twists, its localization sequence has boundary term
the kernel of the sum map $H^0(D)(-1)\to H^2(E)$, of dimension one.
That is a conditional evaluation of an existing theory.

The outstanding actual analytic problem is concrete. Construct the corrected
fixed-boundary point traces and a two-sided Green/Dirac parametrix with
uniform overlap graph estimates. Prove the needed resolvent regularity and
exclude additional modes. Global compactness and the positive fixed-Dirichlet
gap are established above, including both endpoint corners. Specify the spin domains,
Pin convention and phase-line transmission data before claiming an actual
trace interface. The earlier proof labels of v12 do not resolve these
obligations; the domain counterexample shows why they must be revisited.
The source geometry, capacity theorem and independent spectral results remain
available while this analytic work continues.

\appendix
\section{Claim audit, formal coverage and reproducibility}
The immutable historical authority is the 43-page final v12 PDF with SHA-256
\begin{quote}\small
\nolinkurl{4ebb8f0a7d6e88a6c5cd97f65d224b4e938867cabad54f057d3ceafb24167d5a}
\end{quote}
Physical and printed page numbers agree; Section 7 is on pages 14--16,
not 13--15. The full machine-readable and human-readable claim maps pin
66 named blocks, including definitions and remarks, to this PDF. Abstract,
conclusion, figure and historical status-table assertions inherit those
audited results, not their former ``proven'' labels.

The accompanying Lean project uses Lean/Mathlib 4.33.1 with the unchanged
RuledV5 and StokesV5 source dependencies. It proves actual atlas/rank and
small-circle statements through those identical maps, a genuine polar fold
through existing smooth inverse charts, actual Mellin interval integrals,
actual Gram entry integrals and eigenvalue estimates, and selected finite
boundary identities. The separate companion lists precisely what is encoded.
Capacity integrals, Hilbert operator domains, cohomology, sheaf theory and
the actual Green/Dirac parametrices are not certified by that finite inventory.
External Whitney recognition, Milnor comparison, Gysin/MHS and abstract
Krein/buckling results remain cited dependencies with explicit hypotheses.

Independent exact replay rebuilds the actual derivative determinants and
Euclidean jet invariants, the polar remainder correction, weighted adjoint
coefficient, model logarithmic residual and counterexample calculations.
It supplies neither an infinite-dimensional parametrix nor an unbounded
operator-domain certificate. The two registered ancestor figures were reviewed:
the monodromy schematic is useful but can be mistaken for a uniform clock;
the golden-rectangle image is correct background but adds no proof. The new
source diagram and comparison table instead guide the corrected statements.
All ancestor files and historical sealed artifacts remain unchanged.

\begin{thebibliography}{99}
\bibitem{ruled} A. Hayami, \emph{The Orthogonal-Circle Ruled Surface}, corrected v5,
2026. \href{https://doi.org/10.5281/zenodo.23073642}{doi:10.5281/zenodo.23073642}.
Proof companion version 0.03:
\href{https://doi.org/10.5281/zenodo.23073654}{doi:10.5281/zenodo.23073654}.
\bibitem{whitneymetric} M. Hasegawa et al., \emph{Intrinsic properties of surfaces
with singularities}, \href{https://arxiv.org/abs/1409.0281}{arXiv:1409.0281},
Section 4, especially Corollary 4.5.
\bibitem{milnor} J. Milnor, \emph{Singular Points of Complex Hypersurfaces},
Annals of Mathematics Studies 61, Princeton University Press, 1968.
\bibitem{dimca} A. Dimca, \emph{Singularities and Topology of Hypersurfaces},
Springer, 1992.
\bibitem{deligne} P. Deligne, Theorie de Hodge II, \emph{Publ. Math. IHES} 40
(1971), 5--57; III, 44 (1974), 5--77.
\bibitem{peters} C. Peters and J. Steenbrink, \emph{Mixed Hodge Structures},
Springer, 2008. Smooth open curves and localization.
\bibitem{principal} L. F. Martins, K. Saji, S. P. dos Santos and R. Shimada,
\emph{Geometry on principal curvature surfaces},
\href{https://arxiv.org/abs/2607.21796}{arXiv:2607.21796v1},
Proposition 2.1 and Eq.~(2.8).
\bibitem{reed} M. Reed and B. Simon, \emph{Methods of Modern Mathematical
Physics I: Functional Analysis}, Academic Press, 1972.
\bibitem{brezis} H. Brezis, \emph{Functional Analysis, Sobolev Spaces and
Partial Differential Equations}, Springer, 2011. Rellich compactness and
Sobolev traces on bounded Lipschitz domains.
\bibitem{hardy} H. Hedenmalm, P. Lindqvist and K. Seip, A Hilbert space of
Dirichlet series and systems of dilated functions in $L^2(0,1)$,
\emph{Duke Math. J.} 86 (1997), 1--37.
\href{https://arxiv.org/abs/math/9512211}{arXiv:math/9512211}.
\bibitem{point} D. Noja and F. Raso Stoia, \emph{The Dirichlet Laplacian
with a point interaction on unbounded Lipschitz domains},
\href{https://arxiv.org/abs/2607.04349}{arXiv:2607.04349v1},
Proposition 2.1 and Theorem 2.2. Regular-domain comparison only.
\bibitem{pacini} T. Pacini, \emph{Desingularizing isolated conical singularities:
uniform estimates via weighted Sobolev spaces},
\href{https://arxiv.org/abs/1005.3511}{arXiv:1005.3511}. Conifold hypotheses
are not assumed for a Whitney seam.
\bibitem{krein} G. Fucci et al., \emph{The Krein--von Neumann extension revisited},
\href{https://arxiv.org/abs/2102.00685}{arXiv:2102.00685}.
\bibitem{buckling} M. S. Ashbaugh et al., \emph{The Krein--von Neumann
extension and its connection to an abstract buckling problem},
\href{https://arxiv.org/abs/0907.1439}{arXiv:0907.1439}.
\bibitem{stacks} The Stacks Project Authors, Open immersions and extension by
zero, \href{https://stacks.math.columbia.edu/tag/009Z}{Tag 009Z}.
\bibitem{cone} T. Krainer and G. A. Mendoza, \emph{Elliptic complexes of
first-order cone operators: ideal boundary conditions},
\href{https://arxiv.org/abs/1611.06526}{arXiv:1611.06526}. The cone
hypotheses require a separate actual-metric verification.
\bibitem{dirac} C. Bar and W. Ballmann, \emph{Guide to boundary value problems
for Dirac-type operators},
\href{https://arxiv.org/abs/1307.3021}{arXiv:1307.3021}, Section 2.
\bibitem{pin} R. C. Kirby and L. R. Taylor, Pin structures on low-dimensional
manifolds, \emph{Geometry of Low-Dimensional Manifolds 2}, Cambridge
University Press, 1991, 177--242.
\bibitem{aps} M. F. Atiyah, V. K. Patodi and I. M. Singer, Spectral asymmetry
and Riemannian geometry I, \emph{Math. Proc. Cambridge Philos. Soc.} 77
(1975), 43--69. \href{https://doi.org/10.1017/S0305004100049410}{doi:10.1017/S0305004100049410}.
\end{thebibliography}
\end{document}
